% !TeX program = xelatex
% !TeX encoding = UTF-8
% Standalone: compile this file with XeLaTeX.
% To input into a XePersian document, first define \EmergencyContraceptionInNotebook.
% Source: Google Drive notebook/emergency contraception, handwritten pages 1--4.
% pic01.jpg = 20261013_070005694.jpg (page 1)
% pic02.jpg = 20261013_070012787.jpg (page 2)
% pic03.jpg = 20261013_070018998.jpg (page 3)
% pic04.jpg = 20261013_070024030.jpg (page 4)
% A transcription of the source, not an update of its medical recommendations.
\ifdefined\EmergencyContraceptionInNotebook
  \let\ECfinish\relax
\else
  \def\ECfinish{\end{document}}
  \documentclass[12pt,a4paper]{article}
  \usepackage[margin=22mm]{geometry}
  \usepackage{amsmath,amssymb}
  \usepackage{enumitem}
  \usepackage{xepersian}
  \IfFontExistsTF{Amiri}{\settextfont{Amiri}}{\settextfont{DejaVu Sans}}
  \setlatintextfont{DejaVu Serif}
  \setlength{\parindent}{0pt}
  \setlength{\parskip}{5pt}
  \setlength{\emergencystretch}{3em}
  \setlist[itemize]{itemsep=2pt,topsep=4pt}
  \begin{document}
\fi
\providecommand{\ECunclear}[1]{\textbf{[خوانش نامطمئن: #1]}}
\providecommand{\ECmath}[1]{\lr{\(#1\)}}
\providecommand{\ECpage}[1]{\clearpage\section*{صفحهٔ #1}}

\begin{center}
  \Large\bfseries\lr{Emergency Contraception}
\end{center}
\section*{صفحهٔ ۱}
% Printed stationery headings and blank date fields are not medical content.
% Arrows and braces in the handwritten notes are represented by nested lists.

\subsection*{اندیکاسیون‌های \lr{EC}}
\begin{itemize}
  \item پس از رابطهٔ جنسی محافظت‌نشده.
  \item پس از رخ‌دادن حوادث مثل عدم استفاده از پیشگیری.
  \item اگر بیمار داروهایش را \lr{miss} کرده باشد.
\end{itemize}

تخمین ریسک بارداری اغلب دشوار است:
\begin{itemize}
  \item یک نوبت رابطهٔ جنسی محافظت‌نشده در اواسط سیکل: ۲۰–۳۰\% شانس بارداری دارد.
  \item این شانس قبل از روز ۱۰، ۲\% و پس از روز ۱۷، ۳\% است.
\end{itemize}

\subsection*{روش‌های \lr{EC}}
\begin{itemize}
  \item \lr{Cu-IUD}.
  \item \lr{LNG-EC} (لوونورژسترل): خوراکی.
  \item \lr{UPA-EC} (یولیپریستال استات): خوراکی.
\end{itemize}

\subsection*{\lr{Cu-IUD}}
\begin{itemize}
  \item مؤثرترین روش \lr{EC} است.
  \item روش انتخابی است.
  \item تا ۵ روز بعد از \lr{UPSI} قابل استفاده است، یا تا ۵ روز پس از تخمک‌گذاری، اگر روز تخمک‌گذاری مشخص باشد.
\end{itemize}

اگر زنان ظرف ۳ ساعت پس از مصرف دارو استفراغ کنند: دوز دارو تکرار شود یا \lr{IUD} نصب شود.

\subsection*{\lr{UPA}}
\begin{itemize}
  \item \lr{Progesterone-Receptor modulator}.
  \item دوز: ۳۰ میلی‌گرم.
  \item تا ۵ روز پس از \lr{UPSI}.
  \item می‌تواند اوولاسیون را به تأخیر بیندازد.
  \item با برخی داروها که مکانیسم تحریک آنزیم کبدی دارند تداخل دارد.
  \item در بیماران آسمی که تحت درمان \lr{CS} خوراکی هستند توصیه نمی‌شود.
  \item اگر زن شیرده است، نباید تا ۱ هفته شیر بدهد.
  \item پس از مصرف \lr{UPA} باید تا ۵ روز برای شروع \lr{OCP} صبر کرد.
  \item \lr{OCP} کارایی \lr{UPA} را کم می‌کند.
  \item در این ۵ روز و تا زمانی که \lr{OCP} اثر کند (۲ روز برای \lr{POP} و ۷ روز برای ترکیبی)، نباید رابطه داشته باشد.
  \item بارداری فعلی را متوقف نمی‌کند و بر جنین بی‌تأثیر است.
  \item اگر در یک سیکل \lr{UPA} مصرف شده، می‌توان در همان سیکل در نوبت دوم باز هم مصرف نمود (پس از \lr{UPSI} مجدد).
\end{itemize}

\ECpage{۲}
\subsection*{\lr{LNG-EC} — لوونورژسترل}
\begin{itemize}
  \item تا ۷۲ ساعت پس از \lr{UPSI} مؤثر است.
  \item ۶۰–۸۰\% از زنانی که قرار بوده باردار شوند را جلوگیری می‌کند.
  \item تک قرص \lr{1.5 mg} خوراکی.
  \item در بیمارانی که داروهای محرک آنزیم کبدی مصرف می‌کنند، دوز \lr{3 mg} داده می‌شود.
  \item بلافاصله با مصرف آن می‌توان \lr{OCP} را شروع کرد.
  \item تا زمان رسیدن \lr{OCP} به اثربخشی: رابطه نداشته باشد.
  \item اگر در یک سیکل مصرف شود و بیمار باز هم \lr{UPSI} داشته باشد، می‌تواند باز هم \lr{LNG-EC} مصرف کند.
\end{itemize}

\subsection*{\lr{Female Sterilization}}
\begin{itemize}
  \item به صورت لاپاروسکوپیک و از طریق \lr{tubal occlusion}: کلیپس‌ها، تیتانیوم، روی لولهٔ فالوپ.
  \item ریسک بروز خطا و بارداری: ۱ در ۲۰۰ است.
  \item قابل بازگشت است.
  \item از طریق \lr{ligation} لولهٔ فالوپ نیز قابل انجام است، ولی ریسک شکست آن اندکی بالاتر است.
  \item روش \lr{tubal cannulation} از طریق هیستروسکوپی تحت بی‌حسی موضعی انجام می‌شود.
  \item نرخ شکست مشابه \lr{tubal occlusion} است.
  \item برگشت‌ناپذیر است.
  \item ۳ ماه طول می‌کشد تا کاملاً اثربخش باشد.
  \item تأیید آن با \lr{HSG}: امروزه کمتر مورد استفاده واقع می‌شود.
  \item این روش‌ها ریسک \lr{EP} را بالا می‌برند.
  \item ارتباطی با خونریزی‌های شدید قاعدگی (\lr{HMB}) ندارند.
\end{itemize}

\subsection*{وازکتومی}
\begin{itemize}
  \item تحت بی‌حسی موضعی قابل انجام است.
  \item تا حدی برگشت‌پذیر است.
  \item باید آزواسپرمی در \lr{Sperm} تأیید شود: ۲۰ نوبت انزال نیاز است، ۱۲ هفته حداقل.
  \item ریسک کنسر بیضه یا بیماری‌های قلبی را افزایش نمی‌دهد.
  \item ممکن است درد مزمن بیضه ایجاد شود.
  \item سطح تستوسترون، لیبیدو، ارکشن و ارگاسم را تغییر نمی‌کند.
  \item نرخ شکست: ۱ در ۲۰۰۰.
\end{itemize}

\ECpage{۳}
\subsection*{ریسک بارداری به دنبال \lr{UPI}}
\textbf{\lr{Low risk}:}
\begin{itemize}
  \item یک اپیزود \lr{UPI} خارج از دورهٔ پنجرهٔ باروری.
  \item \lr{Failure of a contraceptive method}:
  \begin{itemize}
    \item \lr{miss} کردن بیش از ۱ دوز \lr{OCP} خوراکی.
    \item \lr{miss} کردن ۱ هفته در موعد تزریق دپو یا پچ.
  \end{itemize}
  \item پارگی کاندوم.
\end{itemize}

\textbf{\lr{High risk}:}
\begin{itemize}
  \item چند اپیزود \lr{UPI}.
  \item \lr{UPI} در میانهٔ سیکل (۱۲ روز پیش از زمان مورد انتظار شروع \lr{menses}).
  \item پریودهای نامنظم.
\end{itemize}

\subsection*{انتخاب روش \lr{EC}}
\begin{itemize}
  \item ریسک باردارشدن با \lr{IUD} لوونورژسترل و \lr{IUD} مسی زیر ۱\% است، ولی برای \lr{LNG} خوراکی حدود \lr{2.6\%} است.
  \item در برخی افراد بهتر است از \lr{IUD} استفاده کنیم:
  \begin{itemize}
    \item \ECmath{\mathrm{BMI}\geq30}.
    \item \lr{high risk} برای بارداری.
  \end{itemize}
  \item \lr{IUD} برای نصب نیاز به ویزیت شدن دارد، ولی روش خوراکی نیاز به ویزیت ندارد.
\end{itemize}

\subsection*{بررسی پیش از تجویز}
قبل از نصب \lr{IUD-EC} تست بارداری انجام می‌دهیم:
\begin{itemize}
  \item تست منفی: نصب \lr{IUD}.
  \item پس از نصب \lr{IUD} نیاز به تجویز \lr{OCP} ندارد. ظرف ۳–۴ هفته \lr{review} می‌شود. در صورت پریود نشدن: تست بارداری.
  \item تست مثبت: مشاوره.
\end{itemize}

قبل از تجویز دارو برای \lr{EC} نیاز به تست بارداری نیست، ولی اگر آماده بود انجام می‌دهیم.

\subsection*{داروهای قابل تجویز در \lr{EC}}
\begin{itemize}
  \item \lr{Ulipristal Acetate}:
  \begin{itemize}
    \item مؤثرتر از \lr{LNG}.
    \item تا ۵ روز پس از \lr{UPI} قابل استفاده است.
    \item اگر \ECmath{\mathrm{BMI}\geq30} ارجح است.
    \item بیمار باید حداقل پنج روز صبر کند و سپس \lr{OCP} را شروع کند.
  \end{itemize}
  \item \lr{LNG}:
  \begin{itemize}
    \item تا ۳ روز پس از \lr{UPI} قابل استفاده است.
    \item می‌توان \lr{OCP} را به صورت همزمان آغاز کرد.
  \end{itemize}
\end{itemize}

\ECpage{۴}
\subsection*{دوزها و پیگیری}
\begin{itemize}
  \item دوز \lr{Ulipristal Acetate}: تک‌دوز ۳۰ میلی‌گرم خوراکی.
  \item دوز \lr{LNG}: تک‌دوز \lr{1.5} میلی‌گرم خوراکی.
  \item علی‌رغم شروع همزمان \lr{OCP}، ولی تا ۷ روز کاندوم.
  \item با داروی خوراکی، اگر بیمار تا ۳ هفته بعد پریود نشد: تست بارداری.
\end{itemize}

\subsection*{تداخلات دارویی}
\begin{itemize}
  \item داروهایی که اثربخشی \lr{LNG} را در صورت \ECunclear{دریافت} کاهش می‌دهند: دوز ۲ برابر \lr{LNG} می‌دهیم.
  \begin{itemize}
    \item فنی‌توئین، کاربامازپین، ریفامپین، \lr{efavirenz}.
  \end{itemize}
  \item ریفامپین اثر \lr{UPA} را نیز کاهش می‌دهد.
\end{itemize}

\subsection*{شیردهی و گزینهٔ دیگر}
\begin{itemize}
  \item مادر شیرده باید تا ۲۴ ساعت پس از شروع \lr{UPA} شیر را بدوشد و دور بریزد. سایر \lr{EC}ها منعی در شیردهی ندارند.
  \item تجویز میفپریستون، ۱۰–۲۵ میلی‌گرم، به اندازهٔ \lr{LNG} در \lr{EC} مؤثر است.
\end{itemize}

% The source gives different breastfeeding intervals for UPA on pages 1 and 4;
% both are retained in their original locations rather than silently reconciled.
\ECfinish
